%% ─────────────────────────────────────────────────────────────────────────────
%% Cover letter — Talvin Ackbaraly × Qantev
%% Internship: AI Engineer Intern
%% Même charte que le CV (XCharter, un seul accent, filets fins)
%% Compile: latexmk -pdf cover-letter-qantev.tex
%% ─────────────────────────────────────────────────────────────────────────────
\documentclass[11pt, a4paper]{extarticle}

\usepackage[T1]{fontenc}
\usepackage[utf8]{inputenc}
\usepackage[english]{babel}
\usepackage[top=1.6cm, bottom=1.6cm, left=2.2cm, right=2.2cm]{geometry}
\usepackage{XCharter}
\usepackage[letterspace=80]{microtype}
\usepackage{ragged2e}
\usepackage{xcolor}
\usepackage[hidelinks]{hyperref}

\hypersetup{pdftitle={Cover letter — Talvin Ackbaraly}, pdfauthor={Talvin Ackbaraly}}

%% ── Colour system: identique au CV ──────────────────────────────────────────
\definecolor{accent}{HTML}{2563A0}
\definecolor{ink}{HTML}{1A1A1A}
\definecolor{soft}{HTML}{6E6E6E}
\definecolor{hair}{HTML}{C9C9C9}

\hypersetup{colorlinks=true, urlcolor=accent, linkcolor=accent}
\AtBeginDocument{\color{ink}}

\pagestyle{empty}
\setlength{\parindent}{0pt}
\setlength{\JustifyingParindent}{0pt}
\setlength{\parskip}{8pt}
\linespread{1.08}
\microtypesetup{expansion=false}

\begin{document}

%% ── En-tête : expéditeur ─────────────────────────────────────────────────────
%% Coordonnées seulement : l'école et les liens sont déjà dans le CV
{\color{accent}\fontsize{22}{26}\selectfont\scshape\textls[40]{Talvin Ackbaraly}}\par\vspace{3pt}
{\setlength{\parskip}{0pt}
94270 Le Kremlin-Bicêtre, France\\
\href{tel:+33769389008}{\textcolor{ink}{+33 7 69 38 90 08}}\\
\href{mailto:talvin.ackbaraly@icloud.com}{\textcolor{ink}{talvin.ackbaraly@icloud.com}}\par}
\vspace{-4pt}{\color{hair}\rule{\textwidth}{0.4pt}}\par

%% ── Destinataire et date ─────────────────────────────────────────────────────
\vspace{2pt}
\hfill\begin{tabular}[t]{@{}l@{}}
  \textbf{Qantev}\\
  Engineering team\\
  Paris
\end{tabular}

\today

\vspace{2pt}
\textbf{Application for the AI Engineer Intern position — 6 months from February 2027}

\vspace{2pt}
Dear Qantev team,

\justifying
I am a final-year engineering student at EPITA, specialized in computer vision and deep learning,
and I am applying for your AI Engineer internship. What draws me to Qantev is the gap your offer
describes between research and production. I have worked on both sides, and the hardest part has
always been the step in between: turning a model that performs well in evaluation into a system
people can rely on. In health insurance, where a multi-step pipeline acts on real claims, that step
matters even more.

On the AI side, my final-year project with the National Library of France detects and classifies
illustrations in documents fetched through Gallica's IIIF API. I own the detection stage: training
on 10,000 annotated images, benchmarking models and tuning hyperparameters, up to 94.7\,\% mAP50.
I also like to understand models beyond their APIs: I implemented a U-Net from scratch in PyTorch,
without any pretrained weights. This work taught me that a model is only as good as the way it is measured, which is just as true when
iterating on prompts and agent behavior.

On the engineering side, I build and run software in Python and TypeScript. As a freelancer, I
develop an inventory management application for a business client and operate it in production on
my own: NestJS and Next.js, PostgreSQL, automated backups, error tracking with Sentry and uptime
monitoring. At Studiocanal, I shipped features in an existing Java and Angular codebase used by
more than 100 employees, and set up a GitLab CI/CD pipeline on AWS. I also enjoy tracking down
bottlenecks: profiling a real-time video pipeline with Nsight and porting it to the GPU gave a
24-fold speedup.

Finally, I work every day with coding agents such as Claude Code, Codex and Antigravity. Using
them has shown me how much an agent's output depends on its context, its retrieval and the checks
around it. I would now like to move from using agents to building them, and Qantev's agentic
pipelines and RAG applications are exactly where I want to learn to do that well.

I am available for a six-month internship in Paris from February 2027, and I would be glad to
discuss my application with you in an interview.

\vspace{4pt}
Kind regards,

\vspace{10pt}
\hfill Talvin Ackbaraly\hspace*{1cm}

\end{document}
